\chapter{Partition of Unity}

\noindent\textbf{Intuition:} Recall the Gluing Lemma in topology that if two continuous maps agree on the overlapping (open or closed) part, then we can "glue" them together to get a global continuous map. But it's \textit{not} in general true to "glue" together smooth maps defined on \textit{closed} overlapping part and obtain a smooth result. For example, the non-smooth map $f(x)=|x|$ is obtained by "gluing" the smooth maps $f_+: [0,\infty) \to \mathbb{R}$ and $f_-: (-\infty, 0]\to \mathbb{R}$ given by
\[
f_+(x):= x \text{ and } f_-(x):= -x
\]

\section{Cutoff Functions and Smooth Bump Functions}

\begin{lemma}{Basic Smooth Function Lemma}\\
The function $f:\mathbb{R} \to \mathbb{R}$ defined by
\[
f(t):= \begin{cases}
    e^{-1/t} &,t>0\\
    0&,t\leq 0
\end{cases}
\]
is smooth.
\end{lemma}

\begin{lemma}{Smooth Cutoff Function Lemma}\\
Given any real numbers $r_1<r_2$, there exists a smooth function $h:\mathbb{R} \to \mathbb{R}$, s.t.:
\[
h(t)=1 \text{ for } t\leq r_1, 0<h(t)<1 \text{ for } r_1<t<t_2 \text{ and } h(t)=0 \text{ for } t\geq r_2
\]
\end{lemma}

\begin{proof}
    Let $f$ be the function of the Basic Smooth Function Lemma, and then set
    \[
    h(t):= \frac{f(r_2-t)}{f(r_2-t)+ f(t-r_1)}
    \]
\end{proof}

\begin{lemma}{Smooth Bump Function Lemma}\\
Given any positive real numbers $0<r_1<r_2$, there is a smooth function $H:\mathbb{R}^n \to \mathbb{R}$, s.t.:
\[
H\equiv 1 \text{ on } \overline{B}_{r_1} (0), 0<H(x)<1 \text{ for } x\in B_{r_2} (0) \setminus \overline{B}_{r_1} (0) \text{ and } H\equiv 0 \text{ on } \mathbb{R}^n \setminus B_{r_2} (0)
\]
\end{lemma}

\begin{proof}
    Set $H(x):= h(|x|)$, where $h$ is the smooth cutoff function of the Smooth Cutoff Function Lemma.
\end{proof}

\section{Definition and Existence of Partition of Unity}

\begin{definition}{Support}\\
Let $f:X\to \mathbb{R}^k$ be any real-valued or vector-valued function on a topological space $X$. Then the \textit{support of $f$}, denoted as $\text{supp }f$, is the closure of the set of points where $f$ is nonzero:
\[
\text{supp }f:= \overline{\{ p\in X\mid f(p)\neq 0\}}
\]
Furthermore, we define:
\begin{itemize}
    \item If $\text{supp }f$ is contained in some set $U\subseteq M$, then we say that $f$ is \textit{supported in $U$}.

    \item A function $f$ is said to be \textit{completely supported} if $\text{supp }f$ is a compact set.
\end{itemize}
\end{definition}


\noindent\textbf{Example:} If $H$ is the smooth bump function of the Smooth Bump Function Lemma, then
\[
\text{supp }H= \overline{B}_{r_2} (0)
\]

\begin{definition}{Partition of Unity}\\
Let $M$ be a topological space and let $\mathcal{C}= \{U_\alpha \mid \alpha \in \mathcal{A}\}$ be an arbitrary open cover of $M$. A \textit{partition of unity subordinate to $\mathcal{C}$} is a family $\{\psi_\alpha :M\to \mathbb{R} \mid \alpha \in \mathcal{A}\}$ of continuous functions with the following properties:
\begin{enumerate}
    \item $0\leq \psi_\alpha (p) \leq 1 ,\forall \alpha \in \mathcal{A}, p\in M$.
    \item $\text{supp } \psi_\alpha \subseteq U_\alpha ,\forall \alpha\in \mathcal{A}$.
    \item (\textit{Local Finiteness Condition}). The family of supports $\{\text{supp }\psi_ \alpha\mid \alpha\in \mathcal{A} \}$ is locally finite, i.e.: $\forall p\in M,\exists$ open $U_p\ni p$, s.t.: $U_p$ intersects $\text{supp }\psi_\alpha$ for only finitely many values of $\alpha$.

    \item (\textit{Summation Property}). $\forall x\in M, \sum_{\alpha\in \mathcal{A}} \psi_\alpha(p)=1$.
\end{enumerate}
Furthermore, if $M$ is a smooth manifold with or without boundary, then a family $\{\psi_\alpha :M\to \mathbb{R} \mid \alpha\in \mathcal{A} \}$ of smooth functions with the four properties is called a \textit{smooth partition of unity}.
\end{definition}

\noindent\textbf{Remark:} Now we get a tool, partition of unity, for gluing smooth functions. There are basically two strategies to do so:
\begin{itemize}
    \item (\textit{Direct Gluing}). If we can find an open cover $\{U_\alpha\}$ of $M$ satisfying the conditions stated in coro of Local Behavior-of-Smooth Maps Prop: Gluing Lemma-of-Smooth Maps, then we can just apply it.

    \item (\textit{Soft Blending with Partition of Unity}). If the local defs are not guaranteed to agree, then we have to resort to a partition of unity: Consider the function
    \[
    f_{\text{global}} (p):= \sum_{\alpha \in \mathcal{A}} \psi_\alpha (p) f_\alpha(p)
    \]
    By the property 4 in the def of partition of unity, this function makes a smooth weighted average in the overlapping part. But we have two troubles:
    \begin{itemize}
        \item This strategy may change the value of the original functions, so in usage scenarios, we need to check that the initially desired mathematical properties are inherited.
        \item We need to guarantee the common use of this strategy, say, the existence of partitions of unity:
    \end{itemize}
\end{itemize}


\begin{theorem}{Existence-of-Partitions of Unity-on-Smooth Manifolds Theorem}\\
Let $M$ be a smooth manifold with or without boundary and let $\mathcal{C}= \{U_\alpha \mid \alpha\in \mathcal{A}\}$ be any open cover of $M$. Then there exists a smooth partition of unity subordinate to $\mathcal{C}$.
\end{theorem}


\begin{proof}
    We first deal with the case in which $M$ is a smooth manifold without boundary and the case in which $M$ is a smooth manifold with boundary will be omitted here, sharing the same method and with the core that boundary points are covered by regular coordinate half-balls.
    
    Since each $U_\alpha$ is a smooth manifold, then by the Well-Behaved Basis-for-Smooth Manifolds Prop, it admits a basis $\mathcal{B}_\alpha$ consisting of regular coordinate balls. Then $\mathcal{B}:= \bigcup_{\alpha \in \mathcal{A}} \mathcal{B}_\alpha$ is a basis for $\mathcal{T}_M$. By the Paracompactness-of-Topological Manifolds Theorem, $\mathcal{C}$ has a countable, locally finite refinement $\{B_i\} \subseteq \mathcal{B}$. Then by the Closure-and-Unions-of-Locally Finite Families Lemma.1, the cover $\{\overline{B}_i\}$ is also locally finite.

    $\forall i$, since $B_i$ is a regular coordinate ball in some $U_\alpha$, then by the def of regular coordinate balls, $\exists$ a coordinate ball $B_i' \subseteq U_\alpha$, with a smooth coordinate map $\varphi_i :B_i' \to \mathbb{R}^n$, s.t.: $\overline{B}_i \subseteq B_i'$ and that for some $r_i <r_i'$, we have
    \[
    \varphi_i(B_i)= B_{r_i} (0) ,\varphi_i(\overline{B}_i)= \overline{B}_{r_i} (0) \text{ and } \varphi_i(B_i')=B_{r_i'} (0)
    \]
    For each $i$, let $H_i$ be the smooth bump function in the Smooth Bump Function Lemma with $r_1=r_i$ and $r_2=r_i'$ and then define a function $f_i:M \to \mathbb{R}$ by
    \[
    f_i:= \begin{cases}
        H_i\circ \varphi_i &\text{on }B_i'\\
        0 &\text{on } M\setminus \overline{B}_i
    \end{cases}
    \]
    which is well-defined (since $f_i|_{B_i' \setminus \overline{B}_i} \equiv 0$ for both two defs) and smooth with support $\text{supp }f_i=\overline{B}_i$.

    Consider $f:M\to \mathbb{R}$ given by $p\mapsto \sum_i f_i(p)$. By the local finiteness of the cover $\{\overline{B}_i\}$, the sum $\sum_i f_i(p)$ has only finitely many nonzero terms in an open neighborhood of each point $p\in M$ and hence $f$ is a smooth function. Since each $f_i$ is nonnegative everywhere and positive on $B_i$, and every point of $M$ is in some $B_i$, then $f$ is positive everywhere on $M$. Hence, $\forall i$, we can well-definedly define a function $g_i:M\to \mathbb{R}$ by
    \[
    g_i(x):= \frac{f_i (x)}{f(x)}
    \]
    which is smooth. Clearly by def, we have: $0\leq g_i(p) \leq 1$ and $\sum_i g_i\equiv 1$.

    Now we can define the partition of unity subordinate to $\mathcal{C}$: Since the cover $\{B_i'\}$ is refinement of $\mathcal{C}$, then $\forall i, \exists \alpha_i \in \mathcal{A}$, s.t.: $B_i' \subseteq U_{\alpha_i}$. For each $\alpha\in \mathcal{A}$, define $\psi_\alpha:M \to \mathbb{R}$ by
    \[
    \psi_\alpha:= \sum_{i: \alpha_i=\alpha} g_i
    \]
    where if $\{i \mid \alpha_i=\alpha\}= \emptyset$, then we define $\psi_\alpha:=0$. Then we check:
    \begin{enumerate}
        \item By the Closures-and-Unions-of-Local Finite Families Lemma, we have
        \[
        \text{supp }\psi_\alpha= \overline{\bigcup_{i: \alpha_i=\alpha}B_i} \xlongequal{\text{lemma}} \bigcup_{i: \alpha_i=\alpha}\overline{B}_i \subseteq \bigcup_{i: \alpha_i=\alpha} B_i' \subseteq U_\alpha
        \]

        \item $\forall p\in M,0 \leq \psi_\alpha(p)= \sum_{i: \alpha_i=\alpha} g_i(p) \leq \sum_i g_i(p)=1$.
        \item The local finiteness condition of $\{ \text{supp }\psi_\alpha\}$ inherits from $\{f_i\}$.
        \item $\sum_\alpha \psi_\alpha= \sum_ig_i \equiv1$.
    \end{enumerate}
    Moreover, each $\psi_\alpha$ is smooth, so $\{\psi_\alpha\}_{\alpha \in \mathcal{A}}$ is the desired partition of unity.
\end{proof}


\section{Applications}

\subsection{Extension of Bump Functions}

\begin{definition}{Bump Functions}\\
Let $X$ be a topological space, let $A\subseteq X$ be a closed subset and let $U\subseteq X$ be an open subset containing $A$. A continuous function $\psi: X\to \mathbb{R}$ is called a \textit{bump function for $A$ supported in $U$} if
\[
0\leq \psi\leq 1 \text{ on }X, \psi|_A\equiv 1 \text{ and supp } \psi \subseteq U
\]
\end{definition}

\begin{proposition}{Existence-of-Smooth Bump Functions-on-Smooth Manifolds Proposition}\\
Let $M$ be a smooth manifold with or without boundary. For any closed subset $A\subseteq M$ and any open subset $U\subseteq M$ containing $A$, there is a smooth bump function for $A$ supported in $U$.
\end{proposition}


\begin{proof}
    Note that $\{U_0:= U, U_1:=M\setminus A\}$ is an open cover of $M$. Hence, by the Existence-of-Partitions of Unity Theorem, there is a smooth partition of union $\{\psi_0, \psi_1\}$ subordinate to $\{U_0,U_1\}$. Since $\psi_1 \equiv0$ on $A$, then
    \[
    \psi_0|_A =\sum_i (\psi_i|_A)\equiv 1
    \]
    Therefore, $\psi_0$ is the desired smooth bump function.
\end{proof}

\subsection{Extension of Smooth Functions from Closed Subsets}

\begin{definition}{Smooth Maps on Subsets of Manifolds}\\
Let $M,N$ be smooth manifolds with or without boundary and let $A\subseteq M$ be an arbitrary subset. We say that a map $F:A\to N$ is \textit{smooth on $A$} if it has a smooth extension in an open neighborhood of each point, that is, if $\forall p\in A, \exists$ open $W\subseteq M$ containing $p$, and a smooth map $\widetilde{F}:W\to N$ whose restriction to $W\cap A$ agrees with $F$.
\end{definition}

\begin{lemma}{Exention Lemma for Smooth Functions}\\
Let $M$ be a smooth manifold with or without boundary, let $A\subseteq M$ be a closed subset and $f:A\to \mathbb{R}^k$ be a smooth function. For any open subset $U$ containing $A$, there exists a smooth function $\widetilde{f}:M\to \mathbb{R}^k$ such that $\widetilde{f}|_A=f$ and $\text{supp }\widetilde{f} \subseteq U$.
\end{lemma}

\begin{proof}
    $\forall p\in A$, choose an open neighborhood $W_p$ of $p$ and a smooth function $\widetilde{f}_p:W_p \to \mathbb{R}^k$ that agrees with $f$ on $W_p\cap A$. Since $W_p$ and $U$ are open in $M$, then so is $W_p\cap U$ and hence we may assume $W_p \subseteq U$. Note that $\{W_p \mid p\in A\} \cup \{M\setminus A\}$ is an open cover of $M$. Hence, by the Existence-of-Partitions of Unity Thm, there is a smooth partition of unity $\{\psi_p\mid p\in A\} \cup \{\psi_0\}$ subordinate to this cover, with $\text{supp }\psi_p \subseteq W_p$ and $\text{supp }\psi_0 \subseteq M\setminus A$.

    $\forall p\in A$, since $W_p$ is open in $M$, then the product $\psi_p \widetilde{f}_p$, which is smooth on $W_p$, has a smooth extension to all of $M$, which is zero on $M\setminus \text{supp }\psi_p$. Then we can define $\widetilde{f}: M\to \mathbb{R}^k$ by
    \[
    \widetilde{f} (x):=\sum_{p\in A} \psi_p(x) \widetilde{f}_p(x)
    \]
    By the local finiteness of $\{\text{supp }\psi_p\}_{p\in A}$, this sum has only finite number of nonzero terms in an open neighborhood of any point in $M$. Hence, $\widetilde{f}$ is smooth. $\forall x\in A$, since $\text{supp }\psi_0 \subseteq M\setminus A$, then $\psi_0 (x)=0$, so
    \[
    \widetilde{f}(x):= \sum_{p\in A} \psi_p(x) \widetilde{f}_p (x) \xlongequal{\text{def of } \widetilde{f}_p} \sum_{p\in A} \psi_p(x) f(x)= (\psi_0+ \sum_{p\in A} \psi_p (x)) f(x)=f(x)
    \]
    Therefore, $\widetilde{f}$ is an extension of $f$. For the second statement, it follws from the Closures-and-Unions-od-Local Finite Families Lemma.2 that
    \[
    \text{supp }\widetilde{f}= \overline{\bigcup_{p\in A} \text{supp }\psi_p} =\bigcup_{p\in A} \overline{\text{supp }\psi_p}= \bigcup_{p\in A} \text{supp }\psi_p\subseteq U
    \]
\end{proof}

\noindent\textbf{Remark:} \begin{itemize}
    \item If $A$ is \textit{not} closed, then this lemma may fail: (\textit{Counterexample}). Suppose $M=\mathbb{R},A= (0,1)$. Then the function $f:A\to \mathbb{R}$ given by $x\mapsto \frac{1}{x}$ has no extension.
\item The assumption in this lemma that the codomain of $f$ is $\mathbb{R}^k$, and \textit{not} some other smooth manifold, is need: (\textit{Counterexample}). The identity map $\text{id}_{\mathbb{S}^1}$ on $\mathbb{S}^1$ is smooth, but does \textit{not} have even a continuous extension to a map from $\mathbb{R}^2$ to $\mathbb{S}^1$.
\end{itemize}

\subsection{Exaustion Functions}

\begin{definition}{Exhaustion Functions}\\
Let $X$ be a topological space. An \textit{exhaustion function for $X$} is a continuous function $f:X\to \mathbb{R}$ with the property that the set $f^{-1}((-\infty,c])$ is compact $\forall c\in \mathbb{R}$. Here, we call the set $f^{-1}((-\infty ,c])$ a \textit{sublevel set of $f$}.
\end{definition}

\noindent\textbf{Example:} The functions $f:\mathbb{R}^n \to \mathbb{R}$ and $g:\mathbb{B}^n \to \mathbb{R}$ given by
\[
f(x):= |x|^2 \text{ and } g(x):= \frac{1}{1- |x|^2}
\]
are exhaustion functions.

\noindent\textbf{Remark:} It's clear that if $X$ is compact, then any continuous real-valued function on $X$ is an exhaustion, so we are interesting such functions only noncompact manifolds.

\begin{proposition}{Existence-of-Smooth Exhaustion Functions Proposition}\\
Every smooth manifold with or without boundary $M$ admits a smooth positive exhaustion function.
\end{proposition}

\begin{proof}
    Let $\{V_j\}_{j=1}^\infty$ be any countable open cover of $M$ by precompact open subsets and let $\{\psi_j\}_{j=1}^\infty$ be a smooth partition of unity subordinate this cover. Define $f:M\to \mathbb{R}$ by
    \[
    f(p):= \sum_{j=1}^\infty j\psi_j(p)
    \]
    which is smooth, by the local finiteness, and positive since $f(p) \geq \sum\psi_j (p) =1>0$.

    We claim that $f$ is an exhaustion function $\forall c\in \mathbb{R}$, let $N:= [c]+1> c$. If $p \notin \bigcup_{j=1}^N \overline{V}_j$, then $\psi_j(p)=0$ for $1\leq j\leq N$. Hence, 
    \[
    f(p)= \sum_{j=N+1}^\infty j\psi_j(p) \geq \sum_{j=N+1}^\infty N\psi_j(p)= N\sum_{j=1}^\infty \psi_j(p)=N>c
    \]
    We restate it equivalently that if $f(p)\leq c$, then $p\in \bigcup_{j=1}^N \overline{V}_j$. Therefore, $f^{-1} ((-\infty,c]) \subseteq \bigcup_{j=1}^N \overline{V}_j$. Note that $\bigcup_{j=1}^N \overline{V}_j$ is compact and $f^{-1}((-\infty, c])$ is closed by the continuity of $f$, so $f^{-1} ((-\infty,c])$ is compact.
\end{proof}

\subsection{Level Sets}

\begin{theorem}{Level Sets of Smooth Manifolds Theorem}\\
Let $M$ be a smooth manifold. If $K$ is any closed subset of $M$, then there is a smooth nonnegative function $f:M\to \mathbb{R}$, s.t.: $f^{-1}(0)=K$.
\end{theorem}